Physical activity in office environments is considered a critical factor in maintaining employee health and productivity \cite{PAPAGIANNIDIS2020102027}.
Existing studies assessed office activity levels through surveys and external factors \cite{RONGEN2013406,ALHORR2016369,DISHMAN1998344,de2005effect}.
The present study estimates activity levels from 3D features extracted from point-cloud data.
Rongen et al.\ reported small overall effects of workplace health promotion programs and emphasized the need for high-quality randomized controlled trials \cite{RONGEN2013406}.Al Horr et al.\ confirmed that indoor environmental quality factors significantly affect occupant productivity \cite{ALHORR2016369}.
Dishman et al.\ found limited statistically significant effects from workplace physical-activity interventions and noted low scientific quality in many studies \cite{DISHMAN1998344}.
Nooijen et al.\ identified habitual sitting as a major barrier to reducing sedentary behavior \cite{de2005effect}.
Survey-based and environmental-factor approaches do not capture individual deskwork movements during task execution.
